\section{Introduction}
\label{sec:introduction}

A bank that validates its liquidity stress model against the failure of Silicon Valley Bank meets two numbers that do not fit together. The Basel liquidity coverage ratio assumes that deposits leave at rates between 3\% and 40\% over thirty days, and the supervisory stress test of the European Central Bank compounds a small daily rate over six months so that the median survival period of 103 banks is 176 days under its Adverse shock and 122 days under its Extreme shock \citep{bcbs238,ecblist2019,ecbannex2019}. The bank that failed on 10 March 2023 lost about 85\% of its deposits in two days, and the Federal Reserve's review records outflows of more than 40 billion US dollars on 9 March with another 100 billion expected the next day \citep{bcbs555,fedsvb2023}. Both families of assumptions are forward scenarios: they fix a run-off speed and report what the bank survives. A validator who has to say how far a stated scenario is from failure needs the inverse statement, the speed at which the bank just reaches a stated survival period.

That inverse statement is a reverse stress test. The stress-testing principles of the Basel Committee and the literature on reverse tests treat it as a search for the scenario that produces a stated outcome \citep{bcbs155,glasserman2015,breuercsiszar2013,grundke2018}. For liquidity the outcome is the survival period, the first day on which the net liquidity position turns negative, and the search variable that the forward scenarios already share is the speed of the run-off. We study the simplest scenario family that keeps the structure of a supervisory stress: a multiplier $m$ applied to every daily run-off rate of a base scenario, capped so that no rate exceeds one. The reverse stress asks for the smallest multiplier at which the survival period is at most a target of $T$ days.

Three difficulties make the inverse less trivial than the forward calculation. The survival period is an integer day and a bank that never turns negative has no survival period at all, so the function to invert is a step function with a point at infinity. A target that the base scenario already reaches, a target that lies beyond the horizon of the ladder, and a target that no admissible multiplier reaches have no multiplier, and a solver that returns a number for them returns a wrong one. And a multiplier that is reported has to be reproducible: the survival period at the reported value must be the target or lower, and just below it must be longer.

We propose a solver and a measurement protocol that answer these three points. The proposed method defines the reverse stress as the root of a continuous window function of the multiplier, the minimum of the net liquidity position over the first $T$ days; it rounds the root up to six decimals and verifies the threshold property on the neighbouring decimal; and it refuses a target that is not bracketed, with the reason and without a number. The protocol scores a fast-run base, calibrated only on observed runs dated before 2023, against the 2023 rows of the Basel Committee's table of deposit outflows, with the observation of the failed bank as a labelled anchor and never as a forecast \citep{bcbs555}.

The contributions are as follows.
\begin{itemize}
\item A definition of the reverse liquidity stress as a threshold of a continuous, non-increasing window function, with the cap and the rounding rule that make the answer a reproducible number or an explicit refusal (\cref{sec:proposed-method}).
\item Seven propositions with proofs: the monotonicity of the position in the multiplier, the characterisation of the threshold, the soundness and completeness of the refusal rule, a closed form for one segment, an exact round trip with an explicit sufficient condition, two-sided bounds of a hyperbolic approximation, and an implicit-function sensitivity formula (\cref{sec:theory}).
\item A measurement on \kInst{} instances of a seeded corpus of synthetic banks against an independent oracle, with a scan of a fixed multiplier grid, a scan of a fine grid, Brent's method on the survival period and a bisection without the refusal rule as comparators, and an ablation of each component (\cref{sec:results,sec:ablation}).
\item An application to the three committed bases of case C16, a bank whose balance-sheet aggregates are real and whose deposit behaviour is synthetic: the multipliers, their sensitivity to the case conventions, and a score against the run of March 2023 that is red and that we report as such (\cref{sec:results}).
\end{itemize}

The main results are these. On the corpus the proposed solver agrees with the oracle on every instance: it answers none of the \kRefused{} refused instances and misses none of the \kSolved{} solved ones, and every multiplier it returns is a threshold, whereas a scan of the multiplier grid of the case design answers \kGFalse{}\,\% of the refused instances and overshoots the threshold by a median of \kGErrPc{}\,\% in relative terms. On the bank of case C16 the multiplier that brings the survival period to 30 days is \mAdv{} on the Adverse base and \mExt{} on the Extreme base, and the target is refused on the 2023 fast-run base because its ladder is five days long. The fast-run base reproduces \svbRatio{} of the observed two-day outflow of the failed bank, and the multiplier that would reproduce it is \mSvb{} times the calibrated rate. Three of twelve predictions failed, among them the stability of the Adverse multiplier over the joint band of case conventions.
